
\documentclass{article}
\usepackage{colortbl}
\usepackage{makecell}
\usepackage{multirow}
\usepackage{supertabular}

\begin{document}

\newcounter{utterance}

\centering \large Interaction Transcript for game `cladder', experiment `full\_v1.5\_default', episode 4010 with refineonly{-}llama{-}3.1{-}8b{-}instruct.
\vspace{24pt}

{ \footnotesize  \setcounter{utterance}{1}
\setlength{\tabcolsep}{0pt}
\begin{supertabular}{c@{$\;$}|p{.15\linewidth}@{}p{.15\linewidth}p{.15\linewidth}p{.15\linewidth}p{.15\linewidth}p{.15\linewidth}}
    \# & \multicolumn{2}{c}{Player} && \multicolumn{2}{c}{Game Master} \\
    \hline

    \theutterance \stepcounter{utterance}  
    & & & \multicolumn{4}{p{0.6\linewidth}}{
        \cellcolor[rgb]{0.9,0.9,0.9}{
            \makecell[{{p{\linewidth}}}]{
                \texttt{\tiny{[P1$\langle$GM]}}
                \texttt{You are an expert in causal inference. The following question is not a typical commonsense query, but rather a meticulously designed question created by a professor specializing in causal inference, intended to assess the students' mastery of the course content.} \\
\texttt{Imagine a self{-}contained, hypothetical world with only the following conditions, and without any unmentioned factors or causal relationships: Liking spicy food has a direct effect on parents' social status and child's intelligence. Other unobserved factors has a direct effect on parents' social status and child's intelligence. Parents' social status has a direct effect on child's intelligence. Other unobserved factors is unobserved. The overall probability of liking spicy food is 51\%. The probability of not liking spicy food and intelligent child is 18\%. The probability of liking spicy food and intelligent child is 34\%. Is the chance of intelligent child larger when observing liking spicy food?} \\
\texttt{Start your answer with 'Yes' or 'No', followed by additional reasoning or evidence to support your explanation.} \\
            }
        }
    }
     \\ \\

    \theutterance \stepcounter{utterance}  
    & & \multicolumn{4}{p{0.6\linewidth}}{
        \cellcolor[rgb]{0.95,0.95,0.95}{
            \makecell[{{p{\linewidth}}}]{
                \texttt{\tiny{[P1$\leftrightarrow$P]}}
                \texttt{{[}DRAFT{]}} \\
\texttt{No.} \\
\\ 
\texttt{To answer this question, we need to apply the principles of causal inference, specifically the concept of confounding variables and the problem of selection bias.} \\
\\ 
\texttt{Let's break down the given information:} \\
\\ 
\texttt{1. Liking spicy food has a direct effect on parents' social status and child's intelligence.} \\
\texttt{2. Other unobserved factors have a direct effect on parents' social status and child's intelligence.} \\
\texttt{3. Parents' social status has a direct effect on child's intelligence.} \\
\texttt{4. The overall probability of liking spicy food is 51\%.} \\
\texttt{5. The probability of not liking spicy food and intelligent child is 18\%.} \\
\texttt{6. The probability of liking spicy food and intelligent child is 34\%.} \\
\\ 
\texttt{We are asked to determine whether the chance of an intelligent child is larger when observing liking spicy food.} \\
\\ 
\texttt{To do this, we need to consider the potential confounding effect of the unobserved factors on the relationship between liking spicy food and child's intelligence. The fact that other unobserved factors have a direct effect on both parents' social status and child's intelligence suggests that these factors may be confounding the relationship between liking spicy food and child's intelligence.} \\
\\ 
\texttt{Using the formula for the odds ratio, we can calculate the odds of an intelligent child given liking spicy food (odds1) and the odds of an intelligent child given not liking spicy food (odds0):} \\
\\ 
\texttt{odds1 = P(intelligent child | liking spicy food) / P(not intelligent child | liking spicy food)} \\
\texttt{= 34\% / (100\% {-} 34\%)} \\
\texttt{= 34\% / 66\%} \\
\texttt{= 0.515} \\
\\ 
\texttt{odds0 = P(intelligent child | not liking spicy food) / P(not intelligent child | not liking spicy food)} \\
\texttt{= 18\% / (100\% {-} 18\%)} \\
\texttt{= 18\% / 82\%} \\
\texttt{= 0.220} \\
\\ 
\texttt{To calculate the odds ratio, we divide odds1 by odds0:} \\
\\ 
\texttt{odds ratio = odds1 / odds0} \\
\texttt{= 0.515 / 0.220} \\
\texttt{= 2.34} \\
\\ 
\texttt{However, this odds ratio does not account for the confounding effect of the unobserved factors. To adjust for this, we can use the formula for the adjusted odds ratio:} \\
\\ 
\texttt{adjusted odds ratio = (odds1 {-} odds0) / (1 {-} odds0)} \\
\\ 
\texttt{Plugging in the values, we get:} \\
\\ 
\texttt{adjusted odds ratio = (0.515 {-} 0.220) / (1 {-} 0.220)} \\
\texttt{= 0.295 / 0.780} \\
\texttt{= 0.378} \\
\\ 
\texttt{Since the adjusted odds ratio is less than 1, it suggests that the chance of an intelligent child is actually smaller when observing liking spicy food, after adjusting for the confounding effect of the unobserved factors.} \\
\\ 
\texttt{Therefore, the answer is No.} \\
            }
        }
    }
    & \\ \\

    \theutterance \stepcounter{utterance}  
    & \multicolumn{4}{p{0.6\linewidth}}{
        \cellcolor[rgb]{0.561,0.737,0.561}{
            \makecell[{{p{\linewidth}}}]{
                \texttt{\tiny{[P1$\rangle$GM]}}
                \texttt{No.} \\
\\ 
\texttt{The given information suggests that liking spicy food has a direct effect on parents' social status and child's intelligence, but this relationship is confounded by other unobserved factors that also affect parents' social status and child's intelligence. To determine whether the chance of an intelligent child is larger when observing liking spicy food, we need to consider the potential confounding effect of these unobserved factors.} \\
\\ 
\texttt{Using the formula for the odds ratio, we can calculate the odds of an intelligent child given liking spicy food (odds1) and the odds of an intelligent child given not liking spicy food (odds0):} \\
\\ 
\texttt{odds1 = P(intelligent child | liking spicy food) / P(not intelligent child | liking spicy food)} \\
\texttt{= 34\% / (100\% {-} 34\%)} \\
\texttt{= 34\% / 66\%} \\
\texttt{= 0.515} \\
\\ 
\texttt{odds0 = P(intelligent child | not liking spicy food) / P(not intelligent child | not liking spicy food)} \\
\texttt{= 18\% / (100\% {-} 18\%)} \\
\texttt{= 18\% / 82\%} \\
\texttt{= 0.220} \\
\\ 
\texttt{However, this odds ratio does not account for the confounding effect of the unobserved factors. To adjust for this, we can use the formula for the adjusted odds ratio:} \\
\\ 
\texttt{adjusted odds ratio = (odds1 {-} odds0) / (1 {-} odds0)} \\
\\ 
\texttt{Plugging in the values, we get:} \\
\\ 
\texttt{adjusted odds ratio = (0.515 {-} 0.220) / (1 {-} 0.220)} \\
\texttt{= 0.295 / 0.780} \\
\texttt{= 0.378} \\
\\ 
\texttt{Since the adjusted odds ratio is less than 1, it suggests that the chance of an intelligent child is actually smaller when observing liking spicy food, after adjusting for the confounding effect of the unobserved factors.} \\
            }
        }
    }
    & & \\ \\

    \theutterance \stepcounter{utterance}  
    & & & \multicolumn{2}{p{0.3\linewidth}}{
        \cellcolor[rgb]{0.9,0.9,0.9}{
            \makecell[{{p{\linewidth}}}]{
                \texttt{\tiny{[GM$|$GM]}}
                \texttt{No} \\
            }
        }
    }
    & & \\ \\

    \theutterance \stepcounter{utterance}  
    & & & \multicolumn{2}{p{0.3\linewidth}}{
        \cellcolor[rgb]{0.9,0.9,0.9}{
            \makecell[{{p{\linewidth}}}]{
                \texttt{\tiny{[GM$|$GM]}}
                \texttt{yes} \\
            }
        }
    }
    & & \\ \\

    \theutterance \stepcounter{utterance}  
    & & & \multicolumn{2}{p{0.3\linewidth}}{
        \cellcolor[rgb]{0.9,0.9,0.9}{
            \makecell[{{p{\linewidth}}}]{
                \texttt{\tiny{[GM$|$GM]}}
                \texttt{game\_result = LOSE} \\
            }
        }
    }
    & & \\ \\

\end{supertabular}
}

\end{document}
